% Panel (a) of Fig. 1 (one-column figure): the network on top, the effective-press mechanism below, flowing
% downward. The links of the network end in a dot at the species they act on, not in an arrowhead, so that the
% network is not read as a graph of transitions. Arrowheads are used only in the flow of the mechanism.
% Built by prl_fig1.py, which writes prl_fig1_layout.tex (slot size, cm) and prl_fig1_edges.tex (the links, from
% the matrix A of exp_collinear.npz).
\documentclass[tikz,border=0pt]{standalone}
\usepackage[T1]{fontenc}
\usepackage{helvet}
\renewcommand{\familydefault}{\sfdefault}
\usepackage[helvet]{sfmath}
\usetikzlibrary{arrows.meta,calc}
\definecolor{oigreen}{HTML}{009E73}
\colorlet{linkgrey}{black!38}
\begin{document}
\fontsize{7}{8.4}\selectfont
\input{prl_fig1_layout.tex}%
\begin{tikzpicture}[x=1cm, y=1cm,
  tip/.style={-{Stealth[length=2.3pt 1.6, width=2.0pt 1.6]}},      % flow arrows of the mechanism only
  acts/.style={-{Circle[length=1.9pt 1.0, width=1.9pt 1.0]}},         % species j acts on species i: a dot at i
  sp/.style={circle, draw=black, line width=0.55pt, fill=white, minimum size=0.38cm, inner sep=0pt},
  interaction/.style={draw=linkgrey, bend left=14, shorten >=0.9pt},
  lab/.style={inner sep=1pt, font=\fontsize{7}{8.2}\selectfont},
  ]
\useasboundingbox (0,0) rectangle (\slotW,\slotH);
\newcommand{\gcirc}{\tikz[baseline=-0.6ex]\fill (0,0) circle (1.5pt);}
\newcommand{\gsq}{\tikz[baseline=-0.6ex]\fill (-1.35pt,-1.35pt) rectangle (1.35pt,1.35pt);}
\newcommand{\gtri}{\tikz[baseline=-0.6ex]\fill (0,1.8pt) -- (-1.7pt,-1.2pt) -- (1.7pt,-1.2pt) -- cycle;}
\newcommand{\gdia}{\tikz[baseline=-0.6ex]\draw[line width=0.5pt] (0,1.75pt) -- (1.75pt,0) -- (0,-1.75pt) -- (-1.75pt,0) -- cycle;}
% ------------------------------------------------------------------ network
\def\R{0.80}
\coordinate (c) at ({0.5*\slotW+0.05}, {\slotH-1.18});
\foreach \s/\ang in {3/-90, 1/-18, 5/54, 2/126, 4/198} {
  \node[sp] (n\s) at ($(c)+(\ang:\R)$) {\s};
}
\node[sp, fill=oigreen, text=white, font=\bfseries\fontsize{7}{8.4}\selectfont] at (n3) {3};
\input{prl_fig1_edges.tex}
\draw[black, line width=1.0pt, bend left=14, shorten >=0.9pt, acts] (n1) to (n3);                     % A_31, the changed response
\draw[black, line width=0.85pt, shorten >=0.9pt, acts] (n3.188) to[out=188, in=240, looseness=8] (n3.240);      % A_33 < 0, self-regulation
\draw[black, line width=0.95pt, shorten >=0.9pt] ($(n3.center)+(-45:0.80)$) -- (n3.-45);                    % press on r_3: a plain line (its glyph is a dot already)
\node[lab, anchor=north] at ($(n3.center)+(-45:0.80)+(0.05,-0.02)$) {\gcirc\,press $r_3$};
\node[lab, anchor=east] at ($(n3.center)+(-0.52,-0.28)$) {\gsq\,$A_{33}$};
\node[lab, anchor=west] at ($(n1)!0.5!(n3)+(0.20,-0.10)$) {\gtri\,$A_{31}$};
% ------------------------------------------------------------------ mechanism, flowing downward
\def\yi{2.95}
\node[lab, anchor=west] (in1) at (0.12, \yi+0.18) {\gcirc\,$r_3$};
\node[lab, anchor=west] (in2) at ({0.5*\slotW+0.30}, \yi+0.18) {\gsq\,$A_{33}$};
\node[lab, anchor=west] (in3) at (0.12, \yi-0.18) {\gtri\,$A_{31}$};
\node[lab, anchor=west] (in4) at ({0.5*\slotW+0.30}, \yi-0.18) {\gdia\,mixed};
\coordinate (j) at ({0.5*\slotW}, \yi-0.62);
\foreach \q in {1,3} { \draw[black!70, line width=0.55pt] (in\q.east) to[out=0, in=90] (j); }
\foreach \q in {2,4} { \draw[black!70, line width=0.55pt] (in\q.west) to[out=180, in=90] (j); }
\node[draw, rounded corners=2pt, line width=0.55pt, align=center, inner sep=2.0pt, anchor=north] (sbox)
   at ($(j)+(0,-0.22)$) {effective press $s$};
\draw[line width=0.65pt, tip] (j) -- (sbox.north);
\node[lab, anchor=north] (out) at ($(sbox.south)+(0,-0.30)$) {$\mathbf{x}^{\ast}\!=\mathbf{x}^{\ast}_0+s\,\mathbf{b}_3$};
\draw[line width=0.65pt, tip] (sbox.south) -- (out.north);
\end{tikzpicture}
\end{document}
